\chapter*{Connecting written contracts to surfaces}
\addcontentsline{toc}{chapter}{Introduction: Connecting written contracts to surfaces}

The component atlas makes a portion of the system's visible output available beside its written contracts. Its 42 specimens are generated from actual DOGFOOD story surfaces, with a fixed time, profile, viewport and variant recorded for each capture. The SVG artwork retains the cells' glyphs, positions, colours and supported modifiers; matching vector PDFs are used in the book.

These choices make a specimen traceable. A reader can identify which story and variant produced it and compare the illustration with the contract or implementation that owns the surface. The artwork retains the application's theme even though the surrounding book uses its own palette.

\section*{Read a specimen with its context}

A specimen shows one rendered state. Interaction, resizing, other variants and alternative output modes require their own evidence. Use the story identity and the related chapters to follow those questions into the application. Whole surfaces such as Theme Lab and the BlockLab workbench combine more behaviour than a single component illustration can convey.

The final source inventory connects the book's chapters to their original paths and SHA-256 values. Chapter numbers follow this edition's reading order; source paths and stable record labels preserve identity across editorial rearrangement. The companion manifest maps those records to their LaTeX chapter files.

The atlas and inventory serve complementary purposes: one makes rendered output visible, and the other makes the written source traceable. Together they provide a concrete starting point for checking an explanation against the material from which it was built.
